\documentclass{article}
\usepackage{listings}  % for inline code
\usepackage{hyperref}
\usepackage{CJKutf8}


% Define a simple inline code command
\newcommand{\code}[1]{\texttt{#1}}

\title{CS336 Assignment 1 Writeup}
\author{Legallchaperone}
\date{}

\begin{document}
\begin{CJK}{UTF8}{gbsn}

\maketitle

%% ============================================================
\section*{Problem unicode1: Understanding Unicode}
%% ============================================================

\subsection*{(a)}
\code{chr(0)} returns the null character, whose string representation
is \code{'\textbackslash x00'} (U+0000).

\subsection*{(b)}
% TODO: 填入 repr vs print 的区别
The \code{repr()} of the null character displays as \code{'\textbackslash x00'},
showing the escape sequence, whereas its printed representation is invisible
(nothing appears on screen).

\subsection*{(c)}
% TODO: 填入 null character 在文本中的行为
The null character is invisible when printed and does not appear as whitespace;
it is silently embedded in the string and may act as a string terminator
in some systems.

%% ============================================================
\section*{Problem unicode2: Unicode Encodings}
%% ============================================================

\subsection*{(a)}
% TODO: UTF-8 vs UTF-16/UTF-32 的优势
The UTF-16/UTF-32 encodes the same character with longer bytes, which is inefficient, compared to UTF-8.

\subsection*{(b)}
% TODO: decode_utf8_bytes_to_str_wrong 为什么错，给一个反例
This is wrong because it can only decode one byte at a time, but some characters may be decoded into multiple bytes. Also, some single byte cannot be decoded back to any character.
One counter example input is \code{"牛".encode("utf-8")}.
\subsection*{(c)}
% TODO: 给一个无法解码的两字节序列，解释原因
\code{b'\xe7\x89'}
Simply not defined.
\section*{Problem train\_bpe: BPE Tokenizer Training}
%% ============================================================
% TODO: 主要是代码实现，无书面回答

%% ============================================================
\section*{Problem train\_bpe\_tinystories}
%% ============================================================

\subsection*{(a)}
% TODO: 训练时间、内存用量、最长 token 是什么

\subsection*{(b)}
% TODO: profiling 结果，哪部分最慢

%% ============================================================
\section*{Problem train\_bpe\_expts\_owt}
%% ============================================================

\subsection*{(a)}
% TODO: OWT 最长 token

\subsection*{(b)}
% TODO: TinyStories vs OWT tokenizer 的比较

%% ============================================================
\section*{Problem tokenizer\_experiments}
%% ============================================================

\subsection*{(a)}
% TODO: 压缩比 bytes/token

\subsection*{(b)}
% TODO: 用 TinyStories tokenizer 处理 OWT 的结果

\subsection*{(c)}
% TODO: throughput 估计，处理 Pile 需要多久

\subsection*{(d)}
% TODO: 为什么 uint16 合适

%% ============================================================
\section*{Problem transformer\_accounting}
%% ============================================================

\subsection*{(a)}
% TODO: GPT-2 XL 参数量和内存

\subsection*{(b)}
% TODO: FLOPs 列表和总计

\subsection*{(c)}
% TODO: 哪部分 FLOPs 最多

\subsection*{(d)}
% TODO: 不同 GPT-2 size 的 FLOPs 分布对比

\subsection*{(e)}
% TODO: context length 增大到 16384 的影响

%% ============================================================
\section*{Problem learning\_rate\_tuning}
%% ============================================================
% TODO: 不同学习率下 loss 的行为

%% ============================================================
\section*{Problem adamwAccounting}
%% ============================================================

\subsection*{(a)}
% TODO: 内存分解公式

\subsection*{(b)}
% TODO: GPT-2 XL 具体数值，最大 batch size

\subsection*{(c)}
% TODO: AdamW 每步 FLOPs

\subsection*{(d)}
% TODO: 训练 GPT-2 XL 需要多少天

%% ============================================================
\section*{Problem learning\_rate / batch\_size\_experiment}
%% ============================================================
% TODO: 学习率 sweep 和 batch size 实验的学习曲线和分析
% (可插入图片: \includegraphics[width=\linewidth]{fig.png})

%% ============================================================
\section*{Problem generate}
%% ============================================================
% TODO: 生成的文本 dump 和流畅度分析

%% ============================================================
\section*{Ablations}
%% ============================================================

\subsection*{layer\_norm\_ablation}
% TODO

\subsection*{pre\_norm\_ablation}
% TODO

\subsection*{no\_pos\_emb}
% TODO

\subsection*{swiglu\_ablation}
% TODO

%% ============================================================
\section*{Problem main\_experiment (OpenWebText)}
%% ============================================================
% TODO

%% ============================================================
\section*{Problem leaderboard}
%% ============================================================
% TODO
\end{CJK}
\end{document}